\subsection{In practice: which group, where}\label{sec:lie:groups-in-practice}

\emph{Attitude} (IMU filters, flight controllers) lives in $\SO$ or $S^3$. \emph{Camera and robot poses} (tf2, visual odometry, bundle adjustment) live in $\SE$, and planar mobile robots in $\SEtwo$. \emph{Monocular SLAM} cannot observe metric scale, so when it closes a loop the correction between two keyframes is a similarity transform in $\mathrm{Sim}(3)$. Strasdat, Montiel and Davison (\emph{Scale Drift-Aware Large Scale Monocular SLAM}, RSS 2010) introduced this, and ORB-SLAM's loop closing optimises over $\mathrm{Sim}(3)$ for the same reason. \emph{Inertial navigation} groups orientation, velocity and position into one element of $SE_2(3)$, the ``extended poses'' of Barrau and Bonnabel. On this group the IMU dynamics have a special (group-affine) structure that the invariant EKF exploits (\S\ref{sec:lie:where-else}).

The libraries mirror this list. The Python bindings of manif (Deray and Solà), the library written alongside the micro-Lie-theory paper, ship modules for $\RR^n$, $\SOtwo$, $\SEtwo$, $\SO$, $\SE$, $SE_2(3)$ and the Galilean group $SGal(3)$ (\texttt{bindings\_rn.cpp}, \dots, \texttt{bindings\_sgal3.cpp} in its repository). Every group exposes the same operations: \texttt{compose}, \texttt{inverse}, \texttt{act}, \texttt{exp}/\texttt{log}, \texttt{rplus}/\texttt{lplus}, \texttt{rminus}/\texttt{lminus}, \texttt{adj}, \texttt{rjac}/\texttt{ljac}. That uniformity is the practical payoff of this chapter. Write an estimator once against those operations, and it works on any group.
